% Options for packages loaded elsewhere
\PassOptionsToPackage{unicode}{hyperref}
\PassOptionsToPackage{hyphens}{url}
\PassOptionsToPackage{dvipsnames,svgnames,x11names}{xcolor}
\documentclass[
  11pt,
  a4paper,
]{article}
\usepackage{xcolor}
\usepackage[margin=24mm]{geometry}
\usepackage{amsmath,amssymb}
\setcounter{secnumdepth}{-\maxdimen} % remove section numbering
\usepackage{iftex}
\ifPDFTeX
  \usepackage[T1]{fontenc}
  \usepackage[utf8]{inputenc}
  \usepackage{textcomp} % provide euro and other symbols
\else % if luatex or xetex
  \usepackage{unicode-math} % this also loads fontspec
  \defaultfontfeatures{Scale=MatchLowercase}
  \defaultfontfeatures[\rmfamily]{Ligatures=TeX,Scale=1}
\fi
\usepackage{lmodern}
\ifPDFTeX\else
  % xetex/luatex font selection
\fi
% Use upquote if available, for straight quotes in verbatim environments
\IfFileExists{upquote.sty}{\usepackage{upquote}}{}
\IfFileExists{microtype.sty}{% use microtype if available
  \usepackage[]{microtype}
  \UseMicrotypeSet[protrusion]{basicmath} % disable protrusion for tt fonts
}{}
\makeatletter
\@ifundefined{KOMAClassName}{% if non-KOMA class
  \IfFileExists{parskip.sty}{%
    \usepackage{parskip}
  }{% else
    \setlength{\parindent}{0pt}
    \setlength{\parskip}{6pt plus 2pt minus 1pt}}
}{% if KOMA class
  \KOMAoptions{parskip=half}}
\makeatother
\usepackage{color}
\usepackage{fancyvrb}
\newcommand{\VerbBar}{|}
\newcommand{\VERB}{\Verb[commandchars=\\\{\}]}
\DefineVerbatimEnvironment{Highlighting}{Verbatim}{commandchars=\\\{\}}
% Add ',fontsize=\small' for more characters per line
\newenvironment{Shaded}{}{}
\newcommand{\AlertTok}[1]{\textcolor[rgb]{1.00,0.00,0.00}{\textbf{#1}}}
\newcommand{\AnnotationTok}[1]{\textcolor[rgb]{0.38,0.63,0.69}{\textbf{\textit{#1}}}}
\newcommand{\AttributeTok}[1]{\textcolor[rgb]{0.49,0.56,0.16}{#1}}
\newcommand{\BaseNTok}[1]{\textcolor[rgb]{0.25,0.63,0.44}{#1}}
\newcommand{\BuiltInTok}[1]{\textcolor[rgb]{0.00,0.50,0.00}{#1}}
\newcommand{\CharTok}[1]{\textcolor[rgb]{0.25,0.44,0.63}{#1}}
\newcommand{\CommentTok}[1]{\textcolor[rgb]{0.38,0.63,0.69}{\textit{#1}}}
\newcommand{\CommentVarTok}[1]{\textcolor[rgb]{0.38,0.63,0.69}{\textbf{\textit{#1}}}}
\newcommand{\ConstantTok}[1]{\textcolor[rgb]{0.53,0.00,0.00}{#1}}
\newcommand{\ControlFlowTok}[1]{\textcolor[rgb]{0.00,0.44,0.13}{\textbf{#1}}}
\newcommand{\DataTypeTok}[1]{\textcolor[rgb]{0.56,0.13,0.00}{#1}}
\newcommand{\DecValTok}[1]{\textcolor[rgb]{0.25,0.63,0.44}{#1}}
\newcommand{\DocumentationTok}[1]{\textcolor[rgb]{0.73,0.13,0.13}{\textit{#1}}}
\newcommand{\ErrorTok}[1]{\textcolor[rgb]{1.00,0.00,0.00}{\textbf{#1}}}
\newcommand{\ExtensionTok}[1]{#1}
\newcommand{\FloatTok}[1]{\textcolor[rgb]{0.25,0.63,0.44}{#1}}
\newcommand{\FunctionTok}[1]{\textcolor[rgb]{0.02,0.16,0.49}{#1}}
\newcommand{\ImportTok}[1]{\textcolor[rgb]{0.00,0.50,0.00}{\textbf{#1}}}
\newcommand{\InformationTok}[1]{\textcolor[rgb]{0.38,0.63,0.69}{\textbf{\textit{#1}}}}
\newcommand{\KeywordTok}[1]{\textcolor[rgb]{0.00,0.44,0.13}{\textbf{#1}}}
\newcommand{\NormalTok}[1]{#1}
\newcommand{\OperatorTok}[1]{\textcolor[rgb]{0.40,0.40,0.40}{#1}}
\newcommand{\OtherTok}[1]{\textcolor[rgb]{0.00,0.44,0.13}{#1}}
\newcommand{\PreprocessorTok}[1]{\textcolor[rgb]{0.74,0.48,0.00}{#1}}
\newcommand{\RegionMarkerTok}[1]{#1}
\newcommand{\SpecialCharTok}[1]{\textcolor[rgb]{0.25,0.44,0.63}{#1}}
\newcommand{\SpecialStringTok}[1]{\textcolor[rgb]{0.73,0.40,0.53}{#1}}
\newcommand{\StringTok}[1]{\textcolor[rgb]{0.25,0.44,0.63}{#1}}
\newcommand{\VariableTok}[1]{\textcolor[rgb]{0.10,0.09,0.49}{#1}}
\newcommand{\VerbatimStringTok}[1]{\textcolor[rgb]{0.25,0.44,0.63}{#1}}
\newcommand{\WarningTok}[1]{\textcolor[rgb]{0.38,0.63,0.69}{\textbf{\textit{#1}}}}
\usepackage{longtable,booktabs,array}
\newcounter{none} % for unnumbered tables
\usepackage{calc} % for calculating minipage widths
% Correct order of tables after \paragraph or \subparagraph
\usepackage{etoolbox}
\makeatletter
\patchcmd\longtable{\par}{\if@noskipsec\mbox{}\fi\par}{}{}
\makeatother
% Allow footnotes in longtable head/foot
\IfFileExists{footnotehyper.sty}{\usepackage{footnotehyper}}{\usepackage{footnote}}
\makesavenoteenv{longtable}
\setlength{\emergencystretch}{3em} % prevent overfull lines
\providecommand{\tightlist}{%
  \setlength{\itemsep}{0pt}\setlength{\parskip}{0pt}}
\usepackage{mathtools}
\allowdisplaybreaks
\setlength{\emergencystretch}{6em}
\DeclareUnicodeCharacter{2084}{\ensuremath{_4}}
\usepackage{bookmark}
\IfFileExists{xurl.sty}{\usepackage{xurl}}{} % add URL line breaks if available
\urlstyle{same}
\hypersetup{
  pdftitle={Effective Fourier certificates for the full (e=4) Polydegree column},
  colorlinks=true,
  linkcolor={blue},
  filecolor={Maroon},
  citecolor={Blue},
  urlcolor={blue},
  pdfcreator={LaTeX via pandoc}}

\title{Effective Fourier certificates for the full (e=4) Polydegree
column}
\author{}
\date{23 August 2026}

\begin{document}
\maketitle

\textbf{Candidate research article, 23 August 2026}\\
\textbf{Status:} computer-assisted theorem candidate; internally
replayed; not externally peer reviewed

\subsection{Abstract}\label{abstract}

Let \(\mathcal G_{\mathbf d}\) denote the polydegree stratum of plane
polynomial automorphisms over \(\mathbb C\). Lewis, Perry and Straub
gave an explicit coefficient-polynomial criterion implying

\[
\mathcal G_{(d+4)}\subseteq\overline{\mathcal G_{(d,5)}}
\]

and verified it for \(2\le d<20\). We prove the containment for every
\(d\ge2\). The proof makes a previously qualitative Fourier-limit
argument effective. After writing \(d=4m+r\), one residue class has an
exact solution at the origin. For the other three classes, exact
rational coefficient bounds and complex interval Newton estimates
certify all \(14{,}985\) cases \(5\le m<5000\). A uniform analytic
envelope proves every \(m\ge5000\). The finite ledger was recomputed
under normal and optimized Python, producing byte-identical 17 MB
certificates, and a separate fail-closed audit checks coverage,
ordering, hashes and explicit outward-rounded interval endpoints. The
uniform certificate uses exact rational arithmetic together with 256-bit
Arb evaluation only for cancellation-sensitive Fourier constants. This
is a computer-assisted proof of one Polydegree-containment column; it is
not a result about the Jacobian or Hessian conjectures.

\subsection{1. Statement and relation to earlier
work}\label{statement-and-relation-to-earlier-work}

For \(e\ge1\), write \(\mathbf x=(x_1,\ldots,x_e)\). If
\(\mathbf a\in\mathbb Z_{\ge0}^e\), put

\[
|\mathbf a|=\sum_{s=1}^e a_s,
\qquad
\mathbf a\cdot\mathbf N=\sum_{s=1}^e s a_s.
\]

Following Lewis--Perry--Straub {[}1{]}, define

\[
g_{n,e}=\frac1{n+1}
\sum_{\mathbf a\cdot\mathbf N=n}(-1)^{|\mathbf a|}
\binom{n+|\mathbf a|}{a_1,\ldots,a_e,n}\mathbf x^{\mathbf a}.
\tag{1.1}
\]

Their Theorem 4 implies the following specialization criterion.

\subsubsection{Published input 1.1
(Lewis--Perry--Straub)}\label{published-input-1.1-lewisperrystraub}

If there is \(P\in\mathbb C^e\) satisfying

\[
g_{d,e}(P)=\cdots=g_{d+e-2,e}(P)=0,
\quad
g_{d+e-1,e}(P)\ne0,
\quad
a_{d,e}(P)\ne0,
\tag{1.2}
\]

where \(a_{d,e}\) is their determinant polynomial, then

\[
\mathcal G_{(d+e)}\subseteq
\overline{\mathcal G_{(d,e+1)}}.
\tag{1.3}
\]

They verified (1.2) for \(e=4\) and \(2\le d<20\) {[}1, Theorem
13(b){]}. Perry's earlier dissertation verifies a different, weaker
intersection statement through \(d=45\) {[}2{]}; it must not be read as
the containment in (1.3).

Our main result is as follows.

\subsubsection{Theorem A}\label{theorem-a}

For every integer \(d\ge2\),

\[
\boxed{\mathcal G_{(d+4)}\subseteq
\overline{\mathcal G_{(d,5)}}.}
\tag{1.4}
\]

The new range is \(d\ge20\). Its proof is computer-assisted only on the
finite interval \(20\le d<20000\); an analytic certificate covers all
larger degrees. The residue class \(d\equiv1\pmod4\) also follows from
Edo's known divisibility theorem {[}4{]}; the new work is in the other
three residue classes.

The literature audit accompanying this article did not locate an earlier
all-\(d\) proof of (1.4), but that is an absence finding rather than a
priority theorem. We therefore make no claim of being first until
specialists have checked the result and its historical positioning.

\subsection{2. Reduction to a smooth common
zero}\label{reduction-to-a-smooth-common-zero}

Lewis--Perry--Straub define polynomials \(\alpha_{i,j,e}\) and

\[
a_{d,e}=\det(\alpha_{i,j,e})_{
0\le i\le e-1,\ d-1\le j\le d+e-2}.
\]

Direct coefficient differentiation gives

\[
\alpha_{i,j,e}=-\frac{\partial g_{j+1,e}}{\partial x_{i+1}}.
\tag{2.1}
\]

Let

\[
D_{d,e}=\left(\frac{\partial g_{d+r,e}}{\partial x_s}\right)_ {
0\le r\le e-1,\ 1\le s\le e}.
\]

Then \(a_{d,e}=(-1)^e\det D_{d,e}\). Let \(J_s\) be the minor of the
first \(e-1\) rows obtained by deleting column \(s\).

\subsubsection{Proposition 2.1 (Jacobian--Euler
identity)}\label{proposition-2.1-jacobianeuler-identity}

In \(\mathbb Z[x_1,\ldots,x_e]\),

\[
s x_s a_{d,e}=
\sum_{r=0}^{e-1}(-1)^{e+r+s-1}(d+r)g_{d+r,e}M_{r,s},
\tag{2.2}
\]

where \(M_{r,s}\) deletes row \(r\) and column \(s\) from \(D_{d,e}\).
On the zero locus of the first \(e-1\) rows,

\[
s x_s a_{d,e}=(-1)^s(d+e-1)g_{d+e-1,e}J_s.
\tag{2.3}
\]

\paragraph{Proof}\label{proof}

The polynomial \(g_{n,e}\) is weighted homogeneous of weight \(n\) for
\(\operatorname{wt}(x_s)=s\), hence

\[
\sum_{s=1}^e sx_s\frac{\partial g_{n,e}}{\partial x_s}=ng_{n,e}.
\]

Replace column \(s\) of \(D_{d,e}\) by the weighted sum of all columns.
Multilinearity kills every summand except the original \(s\)-column.
Expanding the resulting determinant along that column and using
\(a_{d,e}=(-1)^e\det D_{d,e}\) proves (2.2). Equation (2.3) follows
after the first \(e-1\) rows vanish. \(\square\)

For \(e=4\), we work on \(x_4=1\) and use \(s=4\). It therefore suffices
to find a common zero of \(g_{d,4},g_{d+1,4},g_{d+2,4}\) at which their
anchored \(3\times3\) Jacobian and \(g_{d+3,4}\) are nonzero.

\subsection{3. Exact rescaling}\label{exact-rescaling}

Write

\[
d=4m+r,\qquad 0\le r\le3.
\]

For each \(n\in\{d,d+1,d+2,d+3\}\), write \(n=4M+q\), where
\(M\in\{m,m+1\}\) and \(0\le q\le3\). Put

\[
C_{M,q}=\frac{(-1)^M}{4M+q+1}
\frac{(5M+q)!}{(4M+q)!M!}
\tag{3.1}
\]

and define the row used by the verifier,

\[
P^{(m)}_{M,q}(X)=
\frac{g_{4M+q,4}(X_1/m,X_2/m,X_3/m,1)}{C_{M,q}}.
\tag{3.2}
\]

For \(\mathbf a=(a_1,a_2,a_3)\), set

\[
u=a_1+a_2+a_3,\quad
w=a_1+2a_2+3a_3,\quad
t=\frac{w-q}{4},\quad
\ell=u-t.
\]

A term occurs precisely when \(w\equiv q\pmod4\) and \(0\le t\le M\).
Its coefficient before \(X_1^{a_1}X_2^{a_2}X_3^{a_3}/(a_1!a_2!a_3!)\) is
exactly

\[
c^{(m)}_{M,q}(\mathbf a)=
\frac{(-1)^\ell}{m^u}
\prod_{i=1}^{\ell}(5M+q+i)
\prod_{i=0}^{t-1}(M-i).
\tag{3.3}
\]

Choose

\[
\rho=5^{1/4}e^{\pi i/4},
\qquad \rho^4=-5.
\]

Since

\[
4\ell-q=3a_1+2a_2+a_3,
\]

the limiting coefficient is

\[
c^\infty_q(\mathbf a)=(-5)^\ell=\rho^q
\rho^{3a_1+2a_2+a_3}.
\tag{3.4}
\]

The exact finite-to-limit ratio is

\[
R^{(m)}_{M,q}(\mathbf a)=
\prod_{i=1}^{\ell}\left(1+\frac{q+i}{5M}\right)
\prod_{i=0}^{t-1}\left(1-\frac{i}{M}\right)
\left(\frac Mm\right)^u.
\tag{3.5}
\]

All analytic and finite bounds below start from (3.3) or (3.5); no
fitted coefficient model is used.

\subsection{4. Fourier limit and simple
zeros}\label{fourier-limit-and-simple-zeros}

Define

\[
H_q(X)=\frac14\sum_{j=0}^3 i^{-qj}
\exp\!\left(\rho^3i^jX_1+\rho^2i^{2j}X_2+\rho i^{3j}X_3\right),
\tag{4.1}
\]

and \(L_q=\rho^qH_q\). The root-of-unity filter shows that \(L_q\) has
the coefficients (3.4), so the four finite rows converge to the
corresponding \(L_q\).

Put

\[
E_j=\exp\!\left(\rho^3i^jX_1+\rho^2i^{2j}X_2+\rho i^{3j}X_3\right).
\]

For residue \(r\), choose \(C_r\) with \(C_r^4=(-1)^{r-1}\) and impose

\[
E_j=C_ri^{(r-1)j}.
\tag{4.2}
\]

The compatibility equation is \(\prod_jE_j=1\). Choose logarithms
\(\Lambda_{j,r}\) summing to zero and set

\[
X^*_{s,r}=\frac1{4\rho^{4-s}}
\sum_{j=0}^3\Lambda_{j,r}i^{-sj},\qquad s=1,2,3.
\tag{4.3}
\]

We use the following exact phase choices, specified by their logarithms
divided by \(\pi i\):

{\def\LTcaptype{none} % do not increment counter
\begin{longtable}[]{@{}rl@{}}
\toprule\noalign{}
\(r\) & \((\Lambda_{0,r},\ldots,\Lambda_{3,r})/(\pi i)\) \\
\midrule\noalign{}
\endhead
\bottomrule\noalign{}
\endlastfoot
0 & \((1/4,-1/4,-3/4,3/4)\) \\
2 & \((1/4,3/4,-3/4,-1/4)\) \\
3 & \((-1/2,1/2,-1/2,1/2)\) \\
\end{longtable}
}

The residue \(r=1\) is treated exactly in Section 7. Fourier inversion
gives

\[
H_r(X_r^*)=H_{r+1}(X_r^*)=H_{r+2}(X_r^*)=0,
\qquad H_{r-1}(X_r^*)=C_r\ne0.
\tag{4.4}
\]

Moreover,

\[
\frac{\partial H_q}{\partial X_s}=\rho^{4-s}H_{q-s}.
\tag{4.5}
\]

Thus the first-three-row limit Jacobian is diagonal after accounting for
the row normalizations. The diagonal moduli and the fourth-row modulus
are:

{\def\LTcaptype{none} % do not increment counter
\begin{longtable}[]{@{}rlr@{}}
\toprule\noalign{}
\(r\) & first-three-row diagonal moduli & fourth-row modulus \\
\midrule\noalign{}
\endhead
\bottomrule\noalign{}
\endlastfoot
0 & \((|\rho|^3,|\rho|^3,|\rho|^3)\) & \(|\rho|^3\) \\
2 & \((|\rho|^5,|\rho|^5,|\rho|)\) & \(|\rho|\) \\
3 & \((|\rho|^6,|\rho|^2,|\rho|^2)\) & \(|\rho|^2\) \\
\end{longtable}
}

Consequently the inverse row-sum norm is at most

\[
|\rho|^{-1}<\frac{67}{100}.
\tag{4.6}
\]

The fourth-row limit modulus is at least

\[
|\rho|>\frac{149}{100}.
\tag{4.7}
\]

A 256-bit Arb evaluation proves

\[
|X^*_{1,r}|\le\frac13,\qquad
|X^*_{2,r}|\le\frac{71}{100},\qquad
|X^*_{3,r}|\le\frac34
\tag{4.8}
\]

for \(r=0,2,3\). These enclosures are reproduced by the uniform
verifier.

\subsection{\texorpdfstring{5. Uniform effectivity for
\(m\ge5000\)}{5. Uniform effectivity for m\textbackslash ge5000}}\label{uniform-effectivity-for-mge5000}

This section proves an explicit analytic handoff. Let

\[
\Delta^{(m)}_{M,q}=P^{(m)}_{M,q}-L_q.
\]

We bound \(\Delta\) on the closed polydisc

\[
\|X-X_r^*\|_\infty\le\frac14.
\tag{5.1}
\]

By (4.8), its coordinate moduli are at most

\[
\left(\frac7{12},\frac{24}{25},1\right).
\tag{5.2}
\]

\subsubsection{5.1 Head expansion through degree
50}\label{head-expansion-through-degree-50}

Let \(\omega=M-m\in\{0,1\}\). Expanding (3.5) gives

\[
R^{(m)}_{M,q}(\mathbf a)
=1+\frac{A_{q,\omega}(\mathbf a)}m+O(m^{-2}),
\tag{5.3}
\]

where

\[
A_{q,\omega}(\mathbf a)=
\frac{q\ell}{5}+\frac{\ell(\ell+1)}{10}
-\frac{t(t-1)}2+\omega u.
\tag{5.4}
\]

The corresponding exponential differential operator can be evaluated by
the four Fourier modes. More explicitly, for mode \(j\), put

\[
z_s=\rho^{4-s}i^{sj}X_s,
\]

\[
p_t=-\frac q4+\sum_{s=1}^3\frac{s}{4}z_s,
\qquad
d_t=\sum_{s=1}^3\frac{s^2}{16}z_s,
\]

\[
p_\ell=\frac q4+\sum_{s=1}^3\frac{4-s}{4}z_s,
\qquad
d_\ell=\sum_{s=1}^3\frac{(4-s)^2}{16}z_s,
\qquad
p_u=\sum_{s=1}^3z_s.
\]

The multiplier corresponding to (5.4) is

\[
\mathcal A_{q,\omega,j}=
\frac q5p_\ell+
\frac{p_\ell^2+d_\ell+p_\ell}{10}
-\frac{p_t^2+d_t-p_t}{2}+\omega p_u,
\tag{5.4a}
\]

and therefore

\[
G_{q,\omega}(X)=\frac{\rho^q}{4}
\sum_{j=0}^3i^{-qj}\mathcal A_{q,\omega,j}E_j.
\tag{5.4b}
\]

At each \(X_r^*\), 256-bit Arb arithmetic proves that the largest
first-order row defect is

\[
1.884955592153876\ldots<2.
\tag{5.5}
\]

On (5.1), an exact rational triangle calculation gives

\[
\|G\|_\infty\le
\frac{51944278023}{204800000}
=253.6341700\ldots<254.
\tag{5.6}
\]

For the second-order remainder, let

\[
A^+=\frac{q\ell}{5}+\frac{\ell(\ell+1)}{10}
+\frac{t(t-1)}2+\omega u.
\]

The elementary product estimate

\[
\left|\prod_j(1+\delta_j)-1-\sum_j\delta_j\right|
\le\frac12\left(\sum_j|\delta_j|\right)^2
\exp\!\left(\sum_j|\delta_j|\right)
\tag{5.7}
\]

together with the difference between \(1/M\) and \(1/m\) gives the exact
coefficient remainder

\[
\left|R^{(m)}_{M,q}-1-\frac{A_{q,\omega}}m\right|
\le \frac{\frac65(A^+)^2/2+A^+}{m^2}
\tag{5.8}
\]

for \(u\le50\) and \(m\ge5000\). Summing (5.8) exactly over all
admissible head monomials gives

\[
K_{2,\mathrm{ball}}<15374,
\qquad
K_{2,\mathrm{centre}}<1105.
\tag{5.9}
\]

The verifier checks the stronger exact rational values
\(15373.3723513\ldots\) and \(1104.3997874\ldots\).

\subsubsection{5.2 Tail above degree 50}\label{tail-above-degree-50}

On (5.1), the absolute limit level of total degree \(u\) is bounded by

\[
B_u=\frac{15^u}{u!}.
\tag{5.10}
\]

Indeed, after dropping the congruence restriction, the multinomial
theorem bounds the level by

\[
\frac{|\rho|^q}{u!}
\left(|\rho|^3|X_1|+|\rho|^2|X_2|+|\rho||X_3|\right)^u.
\]

The bounds \(|\rho|<3/2\) and (5.2) make this smaller than (5.10) for
\(u\ge51\). For a finite coefficient, the falling-factor product in
(3.5) is at most one, while

\[
\prod_{i=1}^{\ell}\left(1+\frac{q+i}{5M}\right)
\left(\frac Mm\right)^u
\le
\exp\!\left(\frac{(u+3)^2}{2m}\right).
\]

This proves the finite majorant used below.

For \(u\ge51\), the first-order multiplier in (5.4) is at most \(u^2\).
Indeed,

\[
A^+\le \frac{61}{160}u^2+\frac{17}{10}u\le u^2.
\tag{5.11}
\]

Here we used \(q\le3\), \(\ell\le u\), \(t\le3u/4\) and \(\omega\le1\).
The last inequality already holds for \(u\ge3\).

For \(u\le m\), the finite level is bounded by

\[
B_u\exp\!\left(\frac{(u+3)^2}{2m}\right).
\tag{5.12}
\]

The verifier sums the finite, limit and \(m^{-1}\)-scaled first-order
levels exactly for \(51\le u\le400\). Starting at \(u=401\), the limit
and first-order successive ratios are bounded by \(1/20\); the same is
true for the finite majorant while \(u\le m\). Its successive ratio is

\[
\frac{15}{u+1}
\exp\!\left(\frac{2u+7}{2m}\right),
\]

which decreases on this range. The three first-ratio inequalities are
checked in exact rational arithmetic, with a rational Taylor upper bound
for the one exponential. For \(u>m\), use \(u!\ge(u/e)^u\), \(M\le m+1\)
and \(q\le3\). On the polydisc (5.1), put \(R_X=\max_s|X_s|\); (5.2)
gives \(R_X\le1\), hence \(3R_X\le3\). For every \(m\ge5000\) and
\(u\ge5001\), the direct finite level is below \((13/1000)^u\): using
\(e<3\), its level base is strictly smaller than

\[
\frac{45}{u}+\frac9m+\frac{72}{mu}<\frac{13}{1000}.
\]

Consequently the complete omitted contribution is

\[
T=2.482176974351622\ldots\times10^{-6}
<\frac1{400000}.
\tag{5.13}
\]

The exact fraction for \(T\) is recomputed by the verifier; its
canonical numerator/denominator text is bound into the receipt by
SHA-256.

\subsubsection{5.3 Quantitative Newton
lemma}\label{quantitative-newton-lemma}

We use the following standard form of the Newton--Kantorovich theorem.

\subsubsection{Lemma 5.1}\label{lemma-5.1}

Let \(F\) be holomorphic on a convex neighbourhood of \(x_0\) in
\(\mathbb C^n\). Suppose \(DF(x_0)\) is invertible and

\[
\|DF(x_0)^{-1}\|\le\beta,
\qquad
\|F(x_0)\|\le\delta,
\]

and suppose \(DF\) is Lipschitz with constant \(L\) on the ball of
radius \(2\eta\), where \(\eta=\beta\delta\). If

\[
h=\beta L\eta\le\frac12,
\]

then \(F\) has a zero in that ball. The zero lies within

\[
t_* = \frac{1-\sqrt{1-2h}}{\beta L}\le2\eta
\]

of \(x_0\). If \(2\beta L\eta<1\), its Jacobian is invertible.

\paragraph{Proof}\label{proof-1}

This is the Newton--Kantorovich majorant theorem applied to the scalar
majorant \(\beta\delta-t+(\beta L/2)t^2\). Its smaller root is \(t_*\).
The inequality \(t_*\le2\eta\) follows from \(1-\sqrt{1-2h}\le2h\).
Finally,

\[
\|DF(x_0)^{-1}(DF(x)-DF(x_0))\|
\le2\beta L\eta<1
\]

on the root enclosure, so Neumann inversion proves nonsingularity.
\(\square\)

\subsubsection{5.4 Cauchy and Newton
constants}\label{cauchy-and-newton-constants}

Equations (5.5)--(5.13) give, for every row and every \(m\ge5000\),

\[
\epsilon_m:=\sup_{(5.1)}|\Delta^{(m)}_{M,q}|
\le\frac{254}{m}+\frac{15374}{m^2}+\frac1{400000},
\tag{5.14}
\]

and at the first-three-row centre,

\[
\delta_m:=\max|P^{(m)}_{M,q}(X_r^*)|
\le\frac2m+\frac{1105}{m^2}+\frac1{400000}.
\tag{5.15}
\]

At \(m=5000\),

\[
\epsilon_m\le0.05141746,
\qquad
\delta_m\le0.0004467.
\tag{5.16}
\]

The first-derivative Cauchy row-sum coefficient at the centre is \(12\),
for which we use \(15\). On the inner ball of radius \(1/100\), the
distance to the boundary of (5.1) is \(6/25\). The second-derivative
Cauchy row-sum coefficient is

\[
\frac{12}{(6/25)^2}=\frac{625}{3}<300.
\tag{5.17}
\]

The limit Hessian row sum is at most \(188.468<258\), and the fourth-row
gradient is at most \(27.095<37\). All four comparisons are derived and
checked in the executable receipt.

Let \(\beta_0=67/100\). Neumann inversion gives

\[
\beta_m\le\frac{\beta_0}{1-15\beta_0\epsilon_m}.
\tag{5.18}
\]

Let

\[
L_m=258+300\epsilon_m,
\quad
\eta_m=\beta_m\delta_m,
\quad
h_m=\beta_mL_m\eta_m.
\tag{5.19}
\]

At the endpoint \(m=5000\), exact rational evaluation gives

{\def\LTcaptype{none} % do not increment counter
\begin{longtable}[]{@{}lr@{}}
\toprule\noalign{}
predicate & certified bound \\
\midrule\noalign{}
\endhead
\bottomrule\noalign{}
\endlastfoot
\(15\beta_0\epsilon_m\) & \(0.516745473<1\) \\
\(\beta_m\) & \(1.386432952\) \\
\(L_m\) & \(273.425238\) \\
\(h_m\) & \(0.234775241<1/2\) \\
\(2\eta_m\) & \(0.001238640<1/100\) \\
fourth-row margin & \(1.392752889>0\) \\
transported Jacobian margin & \(0.530449518>0\) \\
\end{longtable}
}

The Newton--Kantorovich theorem therefore supplies a zero of the first
three finite rows within \(2\eta_m\) of \(X_r^*\). The last two margins
show that the fourth row and the first-three-row Jacobian remain nonzero
at that zero. Every defect expression decreases with \(m\), so the
endpoint proves all \(m\ge5000\).

\subsubsection{Theorem 5.2 (uniform analytic
branch)}\label{theorem-5.2-uniform-analytic-branch}

For every \(m\ge5000\) and \(r\in\{0,2,3\}\), the first three rows
associated with \(d=4m+r\) have a common zero at which their anchored
Jacobian and the fourth row are nonzero.

\subsection{6. Finite interval
certificates}\label{finite-interval-certificates}

It remains to cover \(5\le m<5000\) in residues \(0,2,3\).

\subsubsection{6.1 Exact coefficient evaluation and rigorous
tails}\label{exact-coefficient-evaluation-and-rigorous-tails}

For \(m\le30\), the checker evaluates every coefficient in (3.3). For
\(m\ge31\), it evaluates all monomials of total degree at most \(50\).
Let \(R=4/5\), let \(U=51\), and put

\[
E_{k,u}=
\frac{(5m+8+u)^u(3R)^u u^k}{m^u u!R^k},
\qquad k=0,1,2.
\tag{6.1}
\]

This bounds respectively the omitted value, gradient row sum and Hessian
row sum at level \(u\). Support gives \(u\le d+3\le4m+6\), so

\[
\frac{u}{5m+8+u}<\frac12.
\]

Hence, with an Arb upper bound for \(e^{1/2}\), the first omitted ratio
is

\[
q_{k,U}=
\frac{3R(5m+9+U)e^{1/2}}{m(U+1)}
\left(\frac{U+1}{U}\right)^k<1.
\tag{6.2}
\]

The right side decreases at later levels, so the complete tail is
bounded by

\[
\frac{E_{k,U}}{1-q_{k,U}}.
\tag{6.3}
\]

For completeness, after removing constants the \(u\)-dependent factors
in the ratio majorant are respectively

\[
\frac{5m+9+u}{u+1},\qquad
\frac{5m+9+u}{u},\qquad
\frac{(5m+9+u)(u+1)}{u^2}
\]

for \(k=0,1,2\). Each is decreasing, which justifies using the first
omitted ratio for the complete supported tail.

\subsubsection{6.2 Interval-Newton gate}\label{interval-newton-gate}

Proposed centres are exact rational decimal points. For \(5\le m\le16\),
they are stored explicitly; for \(m\ge17\), the checker uses one fixed
formula

\[
X_{m,r}^{\mathrm{loc}}=X_r^*+\frac{Y_r}{m},
\tag{6.4}
\]

where the decimal rational vectors \(Y_r\) are part of the hashed
source. Their numerical origin is irrelevant to acceptance.

The checker evaluates each truncated row, all first derivatives and all
second derivatives in 128-bit complex Arb arithmetic. The derivative box
is the coordinate rectangle of real and imaginary radius \(1/100\). Its
maximum complex displacement is therefore \(\sqrt2/100\), and the
checker proves that the whole rectangle lies in \(|X_s|<4/5\), the
domain used by (6.1).

Let \(J_0\) be the truncated centre Jacobian and let
\(\beta_{\rm tr}=\|J_0^{-1}\|_\infty\). If \(T_1\) is the gradient-tail
bound, Neumann inversion gives

\[
\beta=\frac{\beta_{\rm tr}}{1-\beta_{\rm tr}T_1}.
\tag{6.5}
\]

Let \(\delta\) be the first-three-row residual including \(T_0\), and
let \(L\) be the Hessian row-sum bound on the box including \(T_2\).
Each case is accepted only if the explicit outward-rounded endpoints
prove

\[
h=\beta L(\beta\delta)<\frac12,
\qquad
2\beta\delta<\frac1{100},
\tag{6.6}
\]

and if interval transport proves both the fourth row and the Jacobian
nonzero throughout the resulting root enclosure.

\subsubsection{6.3 Coverage and worst
case}\label{coverage-and-worst-case}

The ledger contains exactly

\[
4995\times3=14985
\]

records, ordered lexicographically by \(m=5,\ldots,4999\) and
\(r=0,2,3\). Thirty-six cases use stored centres and \(14949\) use
(6.4). The independent structural auditor reports the following worst
endpoints:

{\def\LTcaptype{none} % do not increment counter
\begin{longtable}[]{@{}lrr@{}}
\toprule\noalign{}
quantity & certified endpoint & case \\
\midrule\noalign{}
\endhead
\bottomrule\noalign{}
\endlastfoot
maximum \(h\) & \(0.465782636<1/2\) & \(m=17,r=2\) \\
maximum root radius & \(0.006955535<1/100\) & \(m=17,r=2\) \\
minimum fourth-row margin & \(1.475688755>0\) & \(m=17,r=2\) \\
minimum Jacobian margin & \(0.068434729>0\) & \(m=17,r=2\) \\
\end{longtable}
}

The normal and optimized runs produce byte-identical ledgers with
SHA-256

\begin{Shaded}
\begin{Highlighting}[]
\NormalTok{a106c2df4bf6b771b0ee3245d6e0f5c95d1a1cb5485e9b03e93d942f77ccc3a1}
\end{Highlighting}
\end{Shaded}

The independent audit receipt has SHA-256

\begin{Shaded}
\begin{Highlighting}[]
\NormalTok{d4702c869d6d735a3736ece0a19c61232e667b8dca6ba0c7277e1b5cd89923b6}
\end{Highlighting}
\end{Shaded}

\subsubsection{Theorem 6.1 (finite
bridge)}\label{theorem-6.1-finite-bridge}

For every \(5\le m<5000\) and \(r\in\{0,2,3\}\), the first three rows
associated with \(d=4m+r\) have a common zero at which their anchored
Jacobian and the fourth row are nonzero.

\subsection{\texorpdfstring{7. Exact residue
\(r=1\)}{7. Exact residue r=1}}\label{exact-residue-r1}

Let \(d=4m+1\) and set \(X=(0,0,0)\). Congruence excludes a constant
term from rows \(q=1,2,3\), so the first three rows vanish. Their only
linear monomials are respectively \(X_1,X_2,X_3\), and their Jacobian is
diagonal:

\[
\operatorname{diag}\left(
-\frac{5m+2}{m},
-\frac{5m+3}{m},
-\frac{5m+4}{m}
\right).
\tag{7.1}
\]

Its determinant is nonzero for every \(m\ge1\). The wrapped fourth row
has \(q=0,M=m+1\) and constant term exactly \(1\). Thus this residue
class is algebraic and needs no interval computation. This
specialization recovers the \(e=4\), \(d-1\equiv0\pmod4\) case of Edo's
divisibility theorem {[}4{]}.

\subsection{8. Proof of Theorem A}\label{proof-of-theorem-a}

Lewis--Perry--Straub prove (1.4) for \(2\le d<20\). For \(d\ge20\),
write \(d=4m+r\), so \(m\ge5\). Section 7 proves the required
specialization when \(r=1\). Theorem 6.1 proves it for \(r=0,2,3\) and
\(m<5000\), while Theorem 5.2 proves it for \(m\ge5000\).

Unscaling gives

\[
P=(X_1/m,X_2/m,X_3/m,1).
\]

All row-normalization factors \(C_{M,q}\) are nonzero. Coordinate
scaling and row normalization multiply the anchored Jacobian by nonzero
factors, so the certified normalized determinant is nonzero exactly when
the anchored minor in Proposition 2.1 is nonzero. Proposition 2.1 then
supplies the full determinant condition in Published input 1.1. Applying
the Lewis--Perry--Straub implication proves (1.4). \(\square\)

\subsection{9. Reproducibility and trust
boundary}\label{reproducibility-and-trust-boundary}

The proof package pins Python 3.14.6, python-flint 0.8.0 and the Arb
precision used by each checker. Its proof-critical parts are:

\begin{enumerate}
\def\labelenumi{\arabic{enumi}.}
\tightlist
\item
  \texttt{e4\_eventual\_envelope.py}, which recomputes the exact
  head/tail bounds, the 256-bit Fourier enclosures and the threshold
  inequalities;
\item
  \texttt{e4\_certify\_prototype.py}, the
  coefficient/evaluation/interval-Newton core;
\item
  \texttt{e4\_finite\_bridge.py}, the ordered finite producer;
\item
  \texttt{verify\_finite\_ledger.py}, the separate structural and
  endpoint auditor;
\item
  \texttt{verify\_normalization.py}, which independently reconstructs
  78,144 admissible normalized coefficients from the multinomial
  definition and checks 1,446,940 rejected support triples;
\item
  the exact locator file; and
\item
  the normal and optimized finite ledgers and receipts.
\end{enumerate}

The proof-critical frozen digests are listed in four 16-character
groups; spaces are typographical only:

\begin{itemize}
\tightlist
\item
  eventual receipt:
  \texttt{747cd56f4a297774\ 0ac0786055d2cdf0\ 7ed9e7b89521ed91\ 0ad736d6eaf5daa9};
\item
  normalization receipt:
  \texttt{b82f1709aeee6e1f\ cc58a8091125c11d\ cb3ff5b0b9095fc9\ 85c2bf516a62d47e};
\item
  finite v2 ledger:
  \texttt{a106c2df4bf6b771\ b0ee3245d6e0f5c9\ 5d1a1cb5485e9b03\ e93d942f77ccc3a1};
\item
  finite v2 production receipt:
  \texttt{4e14868e1531fb85\ 509dc017533b9261\ 21c130d31ef1eb1b\ 41195d5cfddb5b54};
\item
  independent ledger audit:
  \texttt{d4702c869d6d735a\ 3736ece0a19c6123\ 2e667b8dca6ba0c7\ 277e1b5cd89923b6}.
\end{itemize}

The decimal locators are not trusted as solutions. The finite proof
depends on the interval gates in (6.6), exact coefficient formulas, and
rigorous omitted tails. Likewise, the finite ledger is not independent
reproduction: both runs use the same source and arithmetic library. Its
evidential role is exhaustive replay and optimization-sensitivity
detection.

Arb midpoint-radius arithmetic is described in {[}3{]}. A specialist
verifier can replace the implementation while preserving the
mathematical certificate conditions in Sections 5 and 6.

\subsection{10. Limitations and claim
boundary}\label{limitations-and-claim-boundary}

\begin{itemize}
\tightlist
\item
  The finite range is a computer-assisted proof, not a hand enumeration.
\item
  The implication from the coefficient specialization to the stratum
  containment is the published theorem of Lewis--Perry--Straub; it is
  not re-formalized in a proof assistant here.
\item
  The package has been internally replayed, but not independently
  reproduced or externally peer reviewed.
\item
  The search found no earlier all-degree \(e=4\) proof, but this does
  not establish priority.
\item
  The theorem concerns plane polynomial automorphism strata. It neither
  proves nor refutes the two-dimensional Jacobian conjecture, the
  quartic Hessian conjecture in dimension four, or any other open
  Jacobian/Hessian case.
\end{itemize}

\subsection{References}\label{references}

{[}1{]} D. Lewis, K. Perry and A. Straub, ``An algorithmic approach to
the Polydegree Conjecture for plane polynomial automorphisms,''
\emph{Journal of Pure and Applied Algebra} \textbf{223} (2019),
5346--5359. \url{https://doi.org/10.1016/j.jpaa.2019.04.002}.

{[}2{]} K. A. Perry, \emph{Polydegree properties of polynomial
automorphisms}, PhD dissertation, The University of Alabama, 2016.
\url{https://ir.ua.edu/handle/123456789/2587}.

{[}3{]} F. Johansson, ``Arb: efficient arbitrary-precision
midpoint-radius interval arithmetic,'' \emph{IEEE Transactions on
Computers} \textbf{66} (2017), 1281--1292.
\url{https://doi.org/10.1109/TC.2017.2690633}.

{[}4{]} E. Edo, ``Some families of polynomial automorphisms II,''
\emph{Acta Mathematica Vietnamica} \textbf{32} (2007), 155--168.

\subsection{Declarations}\label{declarations}

\textbf{Data and code availability.} The exact frozen successor archive,
manifest, source files, ledgers and receipts accompany the candidate
release.

\textbf{Author responsibility.} The submitting author is responsible for
the mathematical claims, computational evidence and source selection.

\textbf{Funding.} No dedicated external funding is declared.

\textbf{Competing interests.} No competing interests are declared.

\textbf{AI-use disclosure.} AI systems contributed to source discovery,
proof exploration, software construction, replay, review and drafting.
The author must inspect and approve the final submission. Internal AI
review is not external peer review.

\end{document}
